\chapter{The Lebesgue Integral and Expectation}

You have been computing expectations for many chapters now. For a discrete
random variable you summed $\sum_x x\,\P(X=x)$; for a continuous one you
integrated $\int x f(x)\,\dd x$. Two formulas, two worlds. Sometimes you met a
variable that was part discrete and part continuous --- a payout that is zero
with probability $1/2$ and otherwise exponential --- and you patched the two
formulas together by hand. It worked, but it felt like a confession: we have no
single definition of expectation, only a family of recipes.

This chapter fixes that. There \emph{is} one definition of $\E[X]$, valid for
every random variable on every probability space, discrete, continuous, mixed,
or stranger still. It is the \emph{Lebesgue integral}. Once you have it, the sum
and the Riemann integral both fall out as special cases, the convergence
theorems you have been quietly using become real theorems, and the deepest
limit results of probability --- the laws of large numbers, the martingale
theorems --- rest on solid ground.

The story is also one of the great ideas of twentieth-century mathematics, and
it begins with a change of perspective so simple it can be drawn on a napkin.

\section{Why the Riemann Integral Is Not Enough}
\label{ch37:why-lebesgue}

\begin{intuitionbox}
Riemann slices the \emph{domain}: chop the $x$-axis into thin vertical strips
and add up their areas. Lebesgue slices the \emph{range}: chop the $y$-axis into
thin horizontal layers and ask, for each height, \emph{how much $x$ lives above
that height?} Same area, opposite cut. The Lebesgue cut is the one probability
wants, because ``how much $x$ lives there'' is exactly a measure --- a
probability.
\end{intuitionbox}

The familiar Riemann integral, for all its virtues in calculus, buckles under
the demands of probability. Here is why.

\subsection{The cashier's parable}

Imagine you must total a pile of coins. The Riemann cashier sweeps the coins
left to right in a line and adds them one by one: $5 + 25 + 1 + 10 + 25 + \dots$
The Lebesgue cashier instead \emph{sorts} the coins into stacks --- all the
pennies, all the nickels, all the quarters --- counts each stack, and computes
$1\cdot(\text{\# pennies}) + 5\cdot(\text{\# nickels}) + 25\cdot(\text{\#
quarters})$. Both get the right total. But the second method does not care about
the \emph{order} the coins came in, and it generalizes to wildly irregular
piles. That sorting-by-value is the Lebesgue idea: group the points of the
domain by the \emph{value} of the function, then weigh each group.

\subsection{A function Riemann cannot touch}

\begin{example}[The Dirichlet function]
\label{ch37:dirichlet}
Define $D\colon[0,1]\to\R$ by $D(x)=1$ if $x$ is rational and $D(x)=0$ if $x$ is
irrational. Try to Riemann-integrate it. In every subinterval, no matter how
small, there are both rationals and irrationals, so the lower Riemann sum is
always $0$ and the upper sum is always $1$. They never meet: $D$ is
\emph{not Riemann integrable}.

Yet morally the answer is obvious. The rationals in $[0,1]$ are countable, hence
negligible --- they take up no length. The function is $0$ ``almost
everywhere,'' so its integral should be $0$. Lebesgue's theory delivers exactly
this: $\int_0^1 D\,\dd x = 0$. Riemann cannot, because it is forced to resolve
$D$ strip by strip, and every strip is hopelessly mixed.
\end{example}

\begin{whybox}
The deeper failure of Riemann is not about pathological functions --- it is
about \emph{limits}. Probability is the study of limits: laws of large numbers,
central limit theorems, convergence of estimators. You constantly want to say
``the limit of the expectations is the expectation of the limit,''
$\lim_n \E[X_n] = \E[\lim_n X_n]$. With Riemann integrals this exchange can fail
badly, and there is no clean theorem telling you when it is safe. A pointwise
limit of Riemann-integrable functions need not even be Riemann integrable (build
$D$ as a limit of finite indicator functions of rationals and you have an
example). Lebesgue's integral is \emph{built to survive limits} --- that is its
whole reason for existing, and the convergence theorems of
Section~\ref{ch37:convergence} are the payoff.
\end{whybox}

\begin{historybox}
Henri Lebesgue presented his theory in his 1902 thesis,
\emph{Int\'egrale, longueur, aire}. He described the change of viewpoint with
his own merchant's metaphor: a shopkeeper totalling the day's takings can add
the bills in the order they were received (Riemann), or first sort the
banknotes by denomination and count each kind (Lebesgue). Émile Borel had
already built the notion of measure for sets; Lebesgue's leap was to use measure
to integrate \emph{functions}, and to prove that the resulting integral behaves
beautifully under limits. Within three decades Kolmogorov would take Lebesgue's
machinery and, in his 1933 \emph{Grundbegriffe}, declare that a probability is
just a measure of total mass one and an expectation is just a Lebesgue
integral --- the foundation the rest of this book stands on.
\end{historybox}

\begin{figure}[H]
\centering
\begin{tikzpicture}
\begin{axis}[
    width=6.4cm, height=5cm,
    title={Riemann: slice the domain},
    xlabel={$x$}, ylabel={$f(x)$},
    domain=0:3, samples=100,
    ymin=0, ymax=3, xmin=0, xmax=3,
    xtick=\empty, ytick=\empty, axis lines=left]
\addplot[ybar interval, fill=sourceblue!25, draw=sourceblue!70] coordinates
  {(0,0.3) (0.5,0.375) (1.0,0.6) (1.5,0.975) (2.0,1.5) (2.5,2.175) (3.0,3.0)};
\addplot[densityred, very thick, samples=100] {0.3*x*x+0.3};
\end{axis}
\end{tikzpicture}
\hspace{0.5cm}
\begin{tikzpicture}
\begin{axis}[
    width=6.4cm, height=5cm,
    title={Lebesgue: slice the range},
    xlabel={$x$}, ylabel={$f(x)$},
    domain=0:3, samples=100,
    ymin=0, ymax=3, xmin=0, xmax=3,
    xtick=\empty, ytick=\empty, axis lines=left]
\foreach \y in {0.4,0.8,1.2,1.6,2.0,2.4,2.8}{
  \edef\yval{\y}
  \addplot[draw=targetgreen!70, thick] coordinates {(0,\yval) (3,\yval)};
}
\addplot[densityred, very thick, samples=100] {0.3*x*x+0.3};
\end{axis}
\end{tikzpicture}
\caption{Two ways to compute the same area. Left, Riemann partitions the
$x$-axis into vertical strips. Right, Lebesgue partitions the $y$-axis into
horizontal layers and measures the set of $x$ lying above each height. The
right-hand cut is the one that turns ``area'' into ``probability times value.''}
\label{ch37:fig-two-cuts}
\end{figure}

\section{The Building Blocks: Simple Functions}
\label{ch37:simple}

We build the integral in four stages, each one extending the last:
\emph{indicators}, then \emph{simple functions}, then \emph{nonnegative
measurable functions}, then \emph{general (signed) functions}. This is the
ladder essentially every textbook climbs, and it is worth climbing once
yourself so the convergence theorems later feel inevitable rather than magical.

Throughout, fix a measure space $(\Omega,\mathcal F,\mu)$: a set $\Omega$, a
$\sigma$-algebra $\mathcal F$ of ``measurable'' subsets, and a measure $\mu$
assigning sizes to them. For probability we will later take $\mu=\P$ with
$\P(\Omega)=1$, but nothing in the construction needs that.

\subsection{Step 1: indicators}

The atom of the whole theory is the indicator $\ind_A$ of a measurable set
$A\in\mathcal F$, which is $1$ on $A$ and $0$ off it. Its integral is, by
\emph{decree}, the size of $A$:
\[
  \int_\Omega \ind_A \,\dd\mu \;=\; \mu(A).
\]
There is nothing to prove; this is the seed. ``The integral of a flat
unit-height plateau over $A$ is the measure of $A$.'' In probability, where
$\mu=\P$, this reads $\E[\ind_A]=\P(A)$ --- the expectation of an indicator is
the probability of the event, a fact you have used since the very first chapters
on expectation. Now you see it is not a coincidence but the \emph{definition}
restricted to its simplest case.

\subsection{Step 2: simple functions}

\begin{definition}[Simple function]
\label{ch37:def-simple}
A \emph{simple function} is a measurable function $s\colon\Omega\to\R$ that takes
only finitely many distinct values $a_1,\dots,a_n$. Writing
$A_k=\{\omega: s(\omega)=a_k\}$ for the (measurable) level sets, we have the
canonical representation
\[
  s \;=\; \sum_{k=1}^n a_k\,\ind_{A_k},
\]
where the $A_k$ are disjoint and cover $\Omega$.
\end{definition}

A simple function is a finite, irregular staircase --- but a staircase laid out
in the \emph{range}, not the domain. Each $A_k$ can be any measurable set, no
matter how scattered; it need not be an interval. This freedom is precisely what
let Lebesgue tame the Dirichlet function.

\begin{definition}[Integral of a nonnegative simple function]
\label{ch37:def-int-simple}
For a nonnegative simple function $s=\sum_{k=1}^n a_k \ind_{A_k}$ (canonical
form, $a_k\ge 0$), define
\[
  \int_\Omega s \,\dd\mu \;=\; \sum_{k=1}^n a_k\,\mu(A_k).
\]
\end{definition}

Read it aloud: \emph{value times the measure of the set where that value is
taken, summed over values}. That is the sorted-coins total. Every later integral
is a limit of these.

\begin{walkbox}
A subtlety worth pausing on: a simple function can be written in many ways, not
just canonically. Is $\sum a_k\mu(A_k)$ the same number regardless of
representation? It is, and the proof is the everyday fact that overlapping
descriptions can be refined to a common one. If $s=\sum a_k\ind_{A_k}=\sum
b_j\ind_{B_j}$, intersect the pieces into $C_{kj}=A_k\cap B_j$; on each
$C_{kj}$ both representations agree that the value is $a_k=b_j$, and finite
additivity of $\mu$ makes the two sums collapse to the same total. So the
integral is \emph{well defined}. Convince yourself of this; it is the only
``bookkeeping'' lemma in the whole construction, and everything downstream
inherits its linearity and monotonicity from it.
\end{walkbox}

Two properties of the simple integral are immediate from finite additivity and
will propagate up the ladder:
\begin{itemize}
\item \textbf{Linearity:} $\int (\alpha s + \beta t)\,\dd\mu = \alpha\int
  s\,\dd\mu + \beta\int t\,\dd\mu$ for nonnegative simple $s,t$ and
  $\alpha,\beta\ge 0$.
\item \textbf{Monotonicity:} if $s\le t$ pointwise then $\int s\,\dd\mu \le \int
  t\,\dd\mu$.
\end{itemize}

\subsection{Step 3: nonnegative measurable functions}

Here is the master stroke. To integrate an arbitrary nonnegative measurable
function $f\ge 0$, we do not approximate it by staircases from the
\emph{outside} (which is what dooms Riemann); we squeeze it from \emph{below} by
simple functions and take the best lower estimate.

\begin{definition}[Lebesgue integral, nonnegative case]
\label{ch37:def-int-nonneg}
For measurable $f\colon\Omega\to[0,\infty]$,
\[
  \int_\Omega f \,\dd\mu \;=\; \sup\set[\Big]{\, \int_\Omega s\,\dd\mu \;:\; s
  \text{ simple},\ 0 \le s \le f \,}.
\]
The supremum is over \emph{all} simple functions sitting under $f$. It always
exists in $[0,\infty]$ (it may be $+\infty$).
\end{definition}

Why does taking the sup of lower estimates work here when it failed for Riemann?
Because Lebesgue's simple functions are allowed to use \emph{arbitrary
measurable level sets}, so they can hug any measurable $f$ from below as tightly
as we like. The following standard construction makes ``as tightly as we like''
concrete and is the engine behind the Monotone Convergence Theorem.

\begin{theorem}[Approximation by simple functions]
\label{ch37:thm-approx}
For every measurable $f\colon\Omega\to[0,\infty]$ there is an increasing
sequence of nonnegative simple functions $s_1\le s_2\le \cdots$ with
$s_n(\omega)\to f(\omega)$ for every $\omega$.
\end{theorem}

\begin{whybox}
The construction is the ``dyadic staircase.'' At stage $n$, slice the range
$[0,\infty)$ into layers of height $2^{-n}$ up to level $n$, and round $f$
\emph{down} to the nearest layer (capping at $n$):
\[
  s_n \;=\; \sum_{k=0}^{n\,2^n-1} \frac{k}{2^n}\,
  \ind_{\set{\frac{k}{2^n}\,\le\, f\,<\,\frac{k+1}{2^n}}}
  \;+\; n\,\ind_{\set{f\ge n}}.
\]
Each level set is measurable because $f$ is, so each $s_n$ is a legitimate
simple function. As $n$ grows, the layers thin out (height $2^{-n}$) and the cap
rises ($n\to\infty$), so $s_n\uparrow f$ at every point. The picture is a
staircase of horizontal --- not vertical --- treads creeping up under the curve.
This single construction is why ``slice the range'' is not a slogan but an
algorithm.
\end{whybox}

\begin{figure}[H]
\centering
\begin{tikzpicture}
\begin{axis}[
    width=11cm, height=5.4cm,
    title={Dyadic simple-function staircase rising under $f$},
    xlabel={$x$}, ylabel={value},
    domain=0:4, samples=200,
    ymin=0, ymax=2.4, xmin=0, xmax=4,
    legend pos=north west, legend cell align=left,
    axis lines=left]
% target function
\addplot[densityred, very thick] {1.8*(1-exp(-0.9*x))+0.05*x};
\addlegendentry{$f(x)$}
% a coarse staircase under it (rounding down to multiples of 0.4)
\addplot[const plot, sourceblue, thick, samples=40] {0.4*floor((1.8*(1-exp(-0.9*x))+0.05*x)/0.4)};
\addlegendentry{$s_n \le f$ (range sliced)}
\end{axis}
\end{tikzpicture}
\caption{The Lebesgue staircase $s_n$ rounds $f$ \emph{down} to a grid of values
on the $y$-axis. Refining the grid (smaller steps, higher cap) pushes
$s_n\uparrow f$ everywhere. Each tread sits over a measurable level set, which is
all the construction ever requires.}
\label{ch37:fig-staircase}
\end{figure}

\subsection{Step 4: general functions}

Finally, any measurable $f\colon\Omega\to\R$ splits into its positive and
negative parts,
\[
  f^+ = \max(f,0), \qquad f^- = \max(-f,0), \qquad f = f^+ - f^-,
\]
both nonnegative and measurable. Define
\[
  \int_\Omega f\,\dd\mu \;=\; \int_\Omega f^+\,\dd\mu \;-\; \int_\Omega
  f^-\,\dd\mu,
\]
\emph{provided at least one of the two pieces is finite} (otherwise we would be
writing $\infty-\infty$, which is meaningless). When both are finite ---
equivalently when $\int |f|\,\dd\mu = \int f^+\,\dd\mu + \int f^-\,\dd\mu <
\infty$ --- we call $f$ \emph{integrable} and write $f\in L^1(\mu)$.

\begin{formalbox}
The full Lebesgue integral, in one breath:
\[
  \int_\Omega f\,\dd\mu
  \;=\; \underbrace{\sup_{0\le s\le f^+}\sum_k a_k\,\mu(A_k)}_{\text{positive part}}
  \;-\; \underbrace{\sup_{0\le s\le f^-}\sum_k a_k\,\mu(A_k)}_{\text{negative part}} .
\]
Build from indicators ($\int\ind_A=\mu(A)$), sum them into simple functions,
take suprema to reach nonnegative functions, and subtract parts for signed
functions. Linearity and monotonicity, established for simple functions, climb
the ladder rung by rung. That is the entire integral.
\end{formalbox}

\section{Expectation Is a Lebesgue Integral}
\label{ch37:expectation}

Now specialize to a probability space $(\Omega,\mathcal F,\P)$, where the measure
is a probability and $\P(\Omega)=1$. A random variable $X$ is a measurable
function $X\colon\Omega\to\R$. Its \emph{expectation} is, simply and finally,
\[
  \boxed{\;\E[X] \;=\; \int_\Omega X\,\dd\P\;}
\]
the Lebesgue integral of $X$ against the probability measure. There is no longer
a discrete formula and a separate continuous formula --- there is one
definition, and the old formulas are theorems about how to \emph{compute} it.

\begin{happeningbox}
What is really happening when we write $\int_\Omega X\,\dd\P$? We are sorting the
outcomes $\omega$ by the value $X(\omega)$ takes, weighing each group of
outcomes by its probability, and adding up value-times-probability. That is the
sorted-coins cashier again, with ``how many coins'' replaced by ``how much
probability.'' Whether $X$ is discrete, continuous, or mixed is irrelevant to
the \emph{definition}; it only changes how you \emph{evaluate} the resulting
integral. The unification you wanted is built into the architecture, not glued
on afterward.
\end{happeningbox}

\subsection{Recovering the discrete sum}

Let $X$ be discrete, taking values $x_1,x_2,\dots$ with $\P(X=x_i)=p_i$. Then
$X$ is exactly a (possibly countable) simple function: $X=\sum_i x_i
\ind_{A_i}$ with $A_i=\{X=x_i\}$ and $\P(A_i)=p_i$. The definition of the simple
integral gives, directly,
\[
  \E[X] \;=\; \int_\Omega X\,\dd\P \;=\; \sum_i x_i\,\P(A_i) \;=\; \sum_i x_i p_i .
\]
The discrete formula is not a different rule --- it is the simple-function
integral, the very first rung of the ladder.

\subsection{Recovering the continuous integral}

Let $X$ be continuous with density $f$. The change-of-variables theorem (the
``law of the unconscious statistician'' in its measure-theoretic form) says that
to integrate against the \emph{distribution} of $X$ on the line we may use the
density:
\[
  \E[X] \;=\; \int_\Omega X\,\dd\P \;=\; \int_\R x \,\dd\mu_X(x) \;=\; \int_\R x\,
  f(x)\,\dd x ,
\]
where $\mu_X$ is the law of $X$ (the pushforward of $\P$). The middle equality
moves the integral from the abstract space $\Omega$ to the real line; the last
equality holds because $\dd\mu_X = f\,\dd x$. The continuous formula, too, is a
special case.

\begin{aestheticbox}
Look at what one definition bought you. The same symbol
$\int X\,\dd\P$ computes a sum when $X$ is discrete, a Riemann-style density
integral when $X$ is continuous, and \emph{something sensible} when $X$ is
neither. You no longer prove the linearity of expectation twice (once for sums,
once for integrals); you prove it once, for the Lebesgue integral, and it holds
universally. Every theorem about expectation in this book --- linearity, the law
of total expectation, Jensen's inequality, the convergence theorems --- is now a
single statement rather than a discrete version plus a continuous version
stapled together.
\end{aestheticbox}

\subsection{The mixed case, painlessly}

\begin{example}[A payout that is part atom, part density]
\label{ch37:mixed}
Let $X$ be the payout of an option: with probability $\tfrac12$ it expires
worthless ($X=0$), and otherwise it pays an amount that is $\Expo(1)$. There is
no density at $0$ (a point mass sits there), and no point mass on $(0,\infty)$
(a density lives there); see Figure~\ref{ch37:fig-mixed-cdf} for the tell-tale
jump-plus-ramp distribution function. Neither the pure-sum nor the pure-integral
formula applies on its own. But the Lebesgue integral does not blink. Split by the
positive/negative-part machinery, or just integrate against the law:
\[
  \E[X] = \underbrace{0\cdot \tfrac12}_{\text{atom at }0}
        + \underbrace{\tfrac12\int_0^\infty x\,e^{-x}\,\dd x}_{\text{density part}}
        = 0 + \tfrac12\cdot 1 = \tfrac12 .
\]
The decomposition into an atomic piece and a continuous piece is automatic ---
it is just $\E[X\ind_{X=0}] + \E[X\ind_{X>0}]$, two Lebesgue integrals over
complementary events, added by linearity. No special-case formula needed.
\end{example}

\begin{figure}[H]
\centering
\begin{tikzpicture}
\begin{axis}[
    width=11cm, height=5cm,
    title={CDF of a mixed variable: a jump (atom) plus a smooth rise (density)},
    xlabel={$t$}, ylabel={$\P(X\le t)$},
    xmin=-1, xmax=5, ymin=0, ymax=1.05,
    domain=0:5, samples=120,
    legend pos=south east, legend cell align=left, axis lines=left]
% atom-free part of the curve for t<0 is 0
\addplot[densityred, very thick] coordinates {(-1,0)(0,0)};
% jump of height 1/2 at t=0 (drawn as a vertical dashed segment)
\addplot[gray, dashed] coordinates {(0,0)(0,0.5)};
% continuous part: 1/2 + (1/2)(1-e^{-t}) for t>=0
\addplot[densityred, very thick, domain=0:5] {0.5+0.5*(1-exp(-x))};
\addlegendentry{$F_X(t)$}
% mark the open/closed dots at the jump
\addplot[only marks, mark=*, densityred, mark size=1.6pt] coordinates {(0,0.5)};
\addplot[only marks, mark=o, densityred, mark size=1.6pt] coordinates {(0,0)};
\end{axis}
\end{tikzpicture}
\caption{The distribution function of the option payout $X$: an atom of mass
$\tfrac12$ at $0$ shows up as a \emph{jump} of height $\tfrac12$ in the CDF,
while the $\Expo(1)$ part contributes the \emph{smooth} climb from $\tfrac12$
toward $1$. A pure pmf would be all jumps; a pure density would be all smooth
rise. The Lebesgue integral handles both halves with one definition.}
\label{ch37:fig-mixed-cdf}
\end{figure}

\begin{applicationbox}
This single-definition viewpoint is not pedantry --- it is what makes modern
probability \emph{computable in code}. When you write a Monte Carlo estimator
$\frac1N\sum_{i=1}^N g(X_i)$ for $\E[g(X)]$, you are approximating one Lebesgue
integral, and you never ask whether $X$ is discrete, continuous, or mixed; you
just sample and average. In machine learning, the expected loss
$\E_{(x,y)\sim\mathcal D}[\ell(f(x),y)]$ is a Lebesgue integral over a data
distribution $\mathcal D$ that is almost always mixed (discrete labels,
continuous features). Empirical risk minimization replaces the true measure $\P$
by the empirical measure $\hat\P_N=\frac1N\sum\delta_{(x_i,y_i)}$ --- a sum of
point masses --- and minimizes $\int \ell\,\dd\hat\P_N$. The reason you can swap
the true distribution for a pile of data points and still call both
``expectation'' is precisely that both are Lebesgue integrals against a measure.
\end{applicationbox}

\section{The Convergence Theorems}
\label{ch37:convergence}

We promised that Lebesgue's integral was built to survive limits. Here is the
payoff: three theorems telling you exactly when
$\lim_n \int f_n = \int \lim_n f_n$. They are the workhorses behind almost every
limit theorem in probability, and you have been implicitly trusting them for
chapters.

\subsection{Monotone Convergence}

\begin{theorem}[Monotone Convergence Theorem, MCT]
\label{ch37:thm-mct}
Let $0\le f_1\le f_2\le\cdots$ be measurable with $f_n\uparrow f$ pointwise.
Then
\[
  \int_\Omega f_n\,\dd\mu \;\longrightarrow\; \int_\Omega f\,\dd\mu .
\]
The limit may be $+\infty$ on both sides; the equality still holds.
\end{theorem}

\begin{whybox}
MCT is almost a restatement of how we \emph{defined} the integral. We set
$\int f = \sup\{\int s : s \text{ simple}, s\le f\}$, and
Theorem~\ref{ch37:thm-approx} produced an increasing sequence of simple
functions climbing to $f$. MCT says any increasing sequence climbing to $f$ ---
not just the dyadic one --- has integrals that climb to $\int f$. The
``$\le$'' direction is monotonicity (each $f_n\le f$, so $\int f_n\le\int f$,
so the limit is $\le\int f$). The ``$\ge$'' direction is the heart: take any
simple $s\le f$ and any $c\in(0,1)$; the sets $E_n=\{f_n\ge cs\}$ increase to
$\Omega$, and continuity of measure from below gives $\int f_n \ge c\int_{E_n} s
\to c\int s$. Let $c\uparrow 1$, then sup over $s$, and the limit of $\int f_n$
dominates $\int f$. The two inequalities pin the limit to $\int f$.
\end{whybox}

\begin{connectionbox}
MCT is the rigorous engine behind a move you have made since the calculus book:
\emph{integrating a series term by term}. If $g_k\ge 0$ are measurable and you
want $\int \sum_k g_k = \sum_k \int g_k$, apply MCT to the partial sums
$f_n=\sum_{k\le n} g_k$, which increase to $\sum_k g_k$. So the Tonelli-style
interchange ``sum of integrals equals integral of the sum'' for nonnegative
terms is just MCT in disguise. The calculus book swapped $\sum$ and $\int$ for
nice power series and asked you to take it on faith; here is the theorem that
licenses it for \emph{any} nonnegative measurable terms. The same logic, applied
to expectations, proves $\E\big[\sum_k Y_k\big]=\sum_k \E[Y_k]$ for nonnegative
$Y_k$ --- countable linearity of expectation, used constantly in renewal theory
and branching processes.
\end{connectionbox}

\subsection{Fatou's Lemma}

A stepping stone, but a useful one in its own right.

\begin{lemma}[Fatou]
\label{ch37:thm-fatou}
For any measurable $f_n\ge 0$,
\[
  \int_\Omega \liminf_{n} f_n \,\dd\mu \;\le\; \liminf_{n} \int_\Omega
  f_n\,\dd\mu .
\]
\end{lemma}

In words: passing a $\liminf$ inside the integral can only \emph{lose} mass,
never gain it. Mass can escape to infinity (a bump that slides off to the right,
or grows ever taller and thinner), and Fatou says that is the only thing that
can go wrong --- the limit integral cannot exceed what the integrals were doing.
It follows from MCT applied to the increasing sequence $g_n=\inf_{k\ge n}f_k$.

\subsection{Dominated Convergence}

This is the one you will reach for most often, because it handles
\emph{signed} functions and gives a clean equality.

\begin{theorem}[Dominated Convergence Theorem, DCT]
\label{ch37:thm-dct}
Suppose $f_n\to f$ pointwise (almost everywhere) and there is a single
integrable ``dominating'' function $g\in L^1(\mu)$ with $|f_n|\le g$ for all
$n$. Then $f\in L^1$ and
\[
  \int_\Omega f_n\,\dd\mu \;\longrightarrow\; \int_\Omega f\,\dd\mu,
  \qquad\text{indeed}\qquad
  \int_\Omega |f_n - f|\,\dd\mu \to 0 .
\]
\end{theorem}

\begin{whybox}
The proof is a clever double application of Fatou. Since $|f_n|\le g$ we have
$g+f_n\ge 0$ and $g-f_n\ge 0$, both nonnegative, so Fatou applies to each.
Applying it to $g+f_n$ gives $\int(g+f)\le\liminf\int(g+f_n)$, i.e. $\int f \le
\liminf\int f_n$. Applying it to $g-f_n$ gives $\int(g-f)\le\liminf\int(g-f_n)$,
i.e. $\int f \ge \limsup\int f_n$. Sandwiching, $\limsup\int f_n \le \int f \le
\liminf\int f_n$, which forces $\int f_n\to\int f$. The dominating $g$ is the
fence that stops mass from escaping to infinity --- the single hypothesis that
rules out Fatou's strict-inequality disasters.
\end{whybox}

\begin{trapbox}
The domination hypothesis is not decoration --- drop it and the theorem is
false. The classic counterexample is the ``escaping bump.'' On $[0,1]$ with
Lebesgue measure let $f_n = n\,\ind_{(0,1/n)}$: a spike of height $n$ on a base
of width $1/n$. Then $f_n(x)\to 0$ for every $x>0$, so $f=0$ and $\int f = 0$.
But $\int f_n = n\cdot\frac1n = 1$ for every $n$, so $\int f_n \to 1 \ne 0$. The
limit of the integrals is not the integral of the limit. Where did the mass go?
It escaped to infinite height as the base shrank. No integrable $g$ can dominate
all the $f_n$ (you would need $g\ge n$ near $0$ for every $n$, so $g=\infty$
there). DCT correctly refuses to apply, and Fatou's inequality holds
\emph{strictly}: $0 = \int\liminf f_n < \liminf\int f_n = 1$.
\end{trapbox}

\begin{figure}[H]
\centering
\begin{tikzpicture}
\begin{axis}[
    width=11cm, height=5cm,
    title={The escaping bump: $f_n = n\,\ind_{(0,1/n)}$, each with integral $1$},
    xlabel={$x$}, ylabel={$f_n(x)$},
    xmin=0, xmax=1.05, ymin=0, ymax=4.6,
    legend pos=north east, legend cell align=left, axis lines=left]
\addplot[const plot, sourceblue, very thick] coordinates {(0,1)(1,1)(1.05,0)};
\addlegendentry{$n=1$}
\addplot[targetgreen, very thick] coordinates {(0,2)(0.5,2)(0.5,0)(1.05,0)};
\addlegendentry{$n=2$}
\addplot[quotientorange, very thick] coordinates {(0,4)(0.25,4)(0.25,0)(1.05,0)};
\addlegendentry{$n=4$}
\end{axis}
\end{tikzpicture}
\caption{Each bump has area exactly $1$, yet $f_n(x)\to 0$ at every fixed
$x>0$. Mass escapes upward as the base narrows, so $\lim\int f_n = 1 \ne 0 =
\int\lim f_n$. Domination by an integrable $g$ is exactly the fence that
prevents this, and DCT requires it.}
\label{ch37:fig-bump}
\end{figure}

\subsection{Why probability cannot live without them}

\begin{aestheticbox}
These theorems are the silent partners in results you have already seen.
\begin{itemize}
\item \textbf{Linearity of expectation for infinite sums}
  ($\E\sum Y_k = \sum \E Y_k$ for $Y_k\ge0$) is MCT.
\item \textbf{Differentiating under the expectation/integral sign}
  --- computing $\frac{\dd}{\dd\theta}\E[g(X,\theta)] = \E[\partial_\theta
  g(X,\theta)]$ in maximum-likelihood and Fisher-information arguments --- is
  DCT applied to difference quotients dominated by an integrable bound.
\item \textbf{The continuity of characteristic functions} and the proof of the
  Central Limit Theorem via $\E[e^{itX_n}]\to\E[e^{itX}]$ is DCT (the
  integrand $e^{itX_n}$ is bounded by the constant $1$, an integrable dominator
  on a probability space).
\item \textbf{Martingale and renewal limit theorems} repeatedly swap $\lim$ and
  $\E$; each swap is a citation to MCT or DCT.
\end{itemize}
Whenever you have written ``$\E$ of the limit is the limit of $\E$,'' you were
spending one of these theorems. Now you know the bank they are drawn on.
\end{aestheticbox}

\begin{paradoxbox}
Here is a puzzle that feels like a contradiction until the convergence theorems
dissolve it. Consider random variables $X_n$ that equal $n$ with probability
$1/n$ and $0$ otherwise. As $n\to\infty$, $\P(X_n\ne 0)=1/n\to 0$, so $X_n\to 0$
in probability --- the variable is almost surely $0$ in the limit. You might
expect $\E[X_n]\to 0$. But $\E[X_n]=n\cdot\frac1n = 1$ for every $n$: the
expectation stubbornly refuses to vanish.

This is the escaping-bump paradox wearing probabilistic clothes. The
\emph{value} $n$ races to infinity exactly as fast as the \emph{probability}
$1/n$ races to zero, and their product stays pinned at $1$. Convergence in
probability does \emph{not} imply convergence of expectations --- you need a
dominating variable (uniform integrability, the probabilist's name for the DCT
hypothesis) to license the interchange. The paradox is not a flaw in
probability; it is the precise warning that expectation is a global, mass-based
quantity that can be hijacked by rare-but-enormous events. Insurance against
catastrophe, heavy-tailed financial risk, and the St.\ Petersburg game all live
in this gap between ``almost never happens'' and ``contributes to the mean.''
\end{paradoxbox}

\section{Almost Everywhere, and What the Integral Cannot See}
\label{ch37:ae}

One more idea makes the whole theory click: the integral is \emph{blind to sets
of measure zero}. A set $N\in\mathcal F$ is \emph{null} if $\mu(N)=0$. A
property holds \emph{almost everywhere} (a.e.), or in probability \emph{almost
surely} (a.s.), if it holds outside some null set.

\begin{proposition}[Null sets are invisible]
\label{ch37:prop-ae}
If $f=g$ almost everywhere, then $\int f\,\dd\mu = \int g\,\dd\mu$. In
particular, if $f=0$ a.e.\ then $\int f\,\dd\mu = 0$, and conversely if $f\ge 0$
with $\int f\,\dd\mu = 0$ then $f=0$ a.e.
\end{proposition}

This is why the Dirichlet function of Example~\ref{ch37:dirichlet} integrates to
$0$: it differs from the constant $0$ only on the rationals, a null set. It is
also why probabilists say ``almost surely'' so often. A statement like ``the
sample mean converges to $\mu$'' (the Strong Law) is an almost-sure statement:
it can fail, but only on a set of outcomes of probability zero, and the integral
--- the expectation --- never notices that set.

\begin{happeningbox}
The right way to think about $L^1$ is that its objects are not really functions
but \emph{equivalence classes} of functions that agree almost everywhere. Two
random variables that differ only on a probability-zero event are, for every
purpose the integral cares about, the same random variable. This is liberating:
you can redefine $X$ arbitrarily on a null set --- to fix a division by zero, to
make it continuous, to tidy a formula --- without changing a single expectation,
variance, or distribution. The integral sees \emph{mass}, and null sets carry
none.
\end{happeningbox}

\section{The Bridge to Riemann}
\label{ch37:riemann-bridge}

We close the circle by relating the new integral to the old one from the
calculus book, so you never have to choose between them.

\begin{theorem}[Lebesgue dominates Riemann]
\label{ch37:thm-riemann}
If $f\colon[a,b]\to\R$ is bounded and Riemann integrable, then it is Lebesgue
integrable on $[a,b]$ and the two integrals agree:
\[
  \int_a^b f(x)\,\dd x \ (\text{Riemann}) \;=\; \int_{[a,b]} f\,\dd\lambda
  \ (\text{Lebesgue}),
\]
where $\lambda$ is Lebesgue measure. Moreover a bounded $f$ on $[a,b]$ is
Riemann integrable \emph{if and only if} its set of discontinuities is a null
set.
\end{theorem}

\begin{connectionbox}
This is the formal handshake with the companion calculus book. Everything you
computed there --- areas under curves, the Fundamental Theorem of Calculus,
$\int_0^\infty x e^{-x}\dd x = 1$ --- is \emph{also} a Lebesgue integral, with
the identical numerical value. So you lose nothing: every Riemann computation you
know transfers verbatim. What you \emph{gain} is (1) functions Riemann cannot
handle (Dirichlet, and limits of nice functions), (2) the convergence theorems
with no analogue in Riemann's world, and (3) a single integral that works on
abstract spaces $\Omega$, not just on intervals of $\R$. The Riemann integral is
the special case ``$\Omega=[a,b]$, $\mu=$ length, $f$ continuous enough''; the
Lebesgue integral is the genuine article. The second sentence of the theorem ---
Riemann integrability is \emph{exactly} ``continuous almost everywhere'' --- is
the sharpest possible statement of why Riemann fails on $D$: $D$ is discontinuous
at \emph{every} point, the opposite of null.
\end{connectionbox}

\begin{trapbox}
The two theories \emph{disagree} for improper integrals. The function
$f(x)=\frac{\sin x}{x}$ has a convergent improper Riemann integral on
$[0,\infty)$ (the famous $\int_0^\infty \frac{\sin x}{x}\dd x = \frac\pi2$),
because the positive and negative areas cancel as you let the upper limit grow.
But $f$ is \emph{not} Lebesgue integrable there, because
$\int_0^\infty \frac{\abs{\sin x}}{x}\dd x = \infty$ --- and the Lebesgue
integral of a signed function requires \emph{absolute} integrability
($\int\abs{f}<\infty$) so that $\int f^+ - \int f^-$ is not $\infty-\infty$. The
lesson: Lebesgue insists on absolute convergence; conditional ``cancellation''
convergence is a Riemann-only trick. For probability this is no loss --- an
expectation $\E[X]$ that depends on the \emph{order} of summation is not a
meaningful average anyway, and we always demand $\E\abs{X}<\infty$ before
speaking of $\E[X]$.
\end{trapbox}

\section{A Fully Worked Computation}
\label{ch37:worked}

We run the entire machine once, end to end, on a single problem, watching the
four stages and the convergence theorems cooperate.

\begin{example}[Expectation of a geometric variable, from the ground up]
\label{ch37:worked-geom}
Let $X$ take value $k$ with probability $\P(X=k)=2^{-k}$ for $k=1,2,3,\dots$
(a $\Geom(1/2)$ counting trials to first success). We compute $\E[X]$ as a
Lebesgue integral and watch the theory work.

\textbf{As a simple-function limit.} $X=\sum_{k\ge1} k\,\ind_{A_k}$ with
$A_k=\{X=k\}$, $\P(A_k)=2^{-k}$ (Figure~\ref{ch37:fig-geom-pmf}). This is a \emph{countable} simple function, so
it is not literally simple (finitely many values), but its truncations
$X_n=\sum_{k=1}^n k\,\ind_{A_k}$ are, and $X_n\uparrow X$. Each truncation
integrates by the simple-function rule:
\[
  \E[X_n]=\int X_n\,\dd\P=\sum_{k=1}^n k\,2^{-k}.
\]

\textbf{Pass the limit by MCT.} Because $0\le X_1\le X_2\le\cdots\uparrow X$,
the Monotone Convergence Theorem licenses
\[
  \E[X]=\lim_{n\to\infty}\E[X_n]=\sum_{k=1}^\infty k\,2^{-k}.
\]
No hand-waving: the swap of $\lim$ and $\int$ is a theorem here, not a hope.

\textbf{Evaluate the series.} Using $\sum_{k\ge1} k r^{k}=\frac{r}{(1-r)^2}$ at
$r=\tfrac12$,
\[
  \E[X]=\frac{1/2}{(1-1/2)^2}=\frac{1/2}{1/4}=2.
\]
So $\E[X]=2$, matching the $\Geom(p)$ mean $1/p$ with $p=\tfrac12$, but now
\emph{derived} from the integral rather than recited from a table. Every step
was a rung of the ladder or a convergence theorem.
\end{example}

\begin{figure}[H]
\centering
\begin{tikzpicture}
\begin{axis}[
    width=11cm, height=4.8cm,
    title={The geometric pmf $\P(X=k)=2^{-k}$ as a countable simple function},
    xlabel={$k$}, ylabel={$\P(X=k)$},
    ybar, bar width=10pt, xtick={1,...,8},
    xmin=0.4, xmax=8.6, ymin=0, ymax=0.56,
    axis lines=left]
\addplot[fill=distblue!35, draw=distblue!80] coordinates
  {(1,0.5)(2,0.25)(3,0.125)(4,0.0625)(5,0.03125)(6,0.015625)(7,0.0078125)(8,0.00390625)};
\end{axis}
\end{tikzpicture}
\caption{Each bar is a level set $A_k=\{X=k\}$ with measure $\P(A_k)=2^{-k}$.
The truncations $X_n=\sum_{k\le n}k\,\ind_{A_k}$ keep only the first $n$ bars and
rise to $X$; the Monotone Convergence Theorem lets their integrals
$\sum_{k\le n}k\,2^{-k}$ climb to $\E[X]=2$.}
\label{ch37:fig-geom-pmf}
\end{figure}

\begin{labbox}
Re-run the machine yourself on the \emph{tail-sum formula}. For a nonnegative
integer-valued $X$, prove
\[
  \E[X]=\sum_{k=1}^\infty \P(X\ge k)
\]
by writing $X=\sum_{k\ge1}\ind_{\{X\ge k\}}$ (check this identity pointwise: if
$X(\omega)=m$, exactly $m$ of the indicators are $1$), then integrating term by
term. Which convergence theorem licenses the interchange of $\sum$ and $\E$? Now
verify it numerically on the $X$ of Example~\ref{ch37:worked-geom}: here
$\P(X\ge k)=2^{-(k-1)}$, so $\sum_{k\ge1}2^{-(k-1)}=2$, agreeing with the
direct computation. The tail-sum formula is one of the most useful tricks in
applied probability, and it is pure MCT.
\end{labbox}

\section{What to Carry Forward}

\begin{focusbox}
The one-paragraph summary of a hundred pages of measure theory:
\begin{itemize}
\item \textbf{One integral.} $\int f\,\dd\mu$ is built in four rungs ---
  indicators ($\int\ind_A=\mu(A)$), simple functions
  ($\int\sum a_k\ind_{A_k}=\sum a_k\mu(A_k)$), nonnegative functions
  (sup of simple functions underneath), signed functions (positive minus
  negative part). It slices the \emph{range}, not the domain.
\item \textbf{One expectation.} $\E[X]=\int_\Omega X\,\dd\P$. The discrete sum
  and the continuous density-integral are both special cases; mixed variables
  need no special handling.
\item \textbf{Limits are safe.} MCT (increasing nonnegative), Fatou (lower
  bound), and DCT (dominated) tell you exactly when $\lim\int=\int\lim$. The
  escaping bump shows what goes wrong without domination.
\item \textbf{Null sets are invisible.} Almost-everywhere / almost-surely is the
  natural language, and the integral never sees a measure-zero set.
\item \textbf{Riemann lives inside.} Every Riemann integral is a Lebesgue
  integral of the same value; Lebesgue adds the limit theorems and the abstract
  domain.
\end{itemize}
With this in hand, the next chapter can finally define conditional expectation
$\E[X\given\mathcal G]$ rigorously --- as a Lebesgue integral relative to a
sub-$\sigma$-algebra --- and the martingale theory rests on bedrock.
\end{focusbox}

\section{Exercises}
\label{ch37:exercises}

\begin{exercise}[Indicators and probabilities]
\label{ch37:ex-indicator}
Let $A,B$ be events. Using only $\E[\ind_E]=\P(E)$ and linearity of the integral,
prove $\P(A\cup B)=\P(A)+\P(B)-\P(A\cap B)$ by integrating the pointwise
identity $\ind_{A\cup B}=\ind_A+\ind_B-\ind_{A\cap B}$. (This is inclusion--
exclusion as a one-line integral.)
\end{exercise}

\begin{exercise}[Well-definedness of the simple integral]
\label{ch37:ex-welldef}
Suppose a simple function has two representations $s=\sum_k a_k\ind_{A_k}=\sum_j
b_j\ind_{B_j}$ with each family disjoint and covering $\Omega$. Prove
$\sum_k a_k\mu(A_k)=\sum_j b_j\mu(B_j)$ by refining with $C_{kj}=A_k\cap B_j$
and using finite additivity. (This is the only bookkeeping lemma the whole
construction needs.)
\end{exercise}

\begin{exercise}[Dirichlet, made precise]
\label{ch37:ex-dirichlet}
Let $D=\ind_{\Q\cap[0,1]}$. (a) Explain why $D$ is a simple function on the
measure space $([0,1],\mathcal B,\lambda)$. (b) Compute its Lebesgue integral
directly from Definition~\ref{ch37:def-int-simple}, using that $\Q\cap[0,1]$ is
countable and hence $\lambda$-null. (c) Show $D$ is not Riemann integrable by
exhibiting its upper and lower Darboux sums.
\end{exercise}

\begin{exercise}[Tail-sum formula]
\label{ch37:ex-tailsum}
For a nonnegative integer-valued random variable $X$, prove
$\E[X]=\sum_{k=1}^\infty\P(X\ge k)$. State clearly which convergence theorem
justifies exchanging $\E$ and the infinite sum, and verify that its hypotheses
hold. Then apply it to compute $\E[X]$ when $X\sim\Geom(p)$ with
$\P(X=k)=(1-p)^{k-1}p$.
\end{exercise}

\begin{exercise}[The escaping bump, your version]
\label{ch37:ex-bump}
On $([0,1],\lambda)$ construct measurable $f_n\ge 0$ with $f_n\to 0$ pointwise
but $\int f_n\,\dd\lambda = 5$ for every $n$. (a) Why does this not contradict
DCT? Identify precisely which hypothesis fails. (b) Compute both
$\int\liminf f_n$ and $\liminf\int f_n$ and confirm Fatou's inequality is
strict.
\end{exercise}

\begin{exercise}[Monotone vs.\ dominated]
\label{ch37:ex-mctdct}
For each sequence on $([0,\infty),\lambda)$, decide whether MCT, DCT, both, or
neither applies, and find $\lim_n\int f_n$.
(a) $f_n(x)=\ind_{[0,n]}(x)\,(1-x/n)_+$.
(b) $f_n(x)=\frac{1}{n}\ind_{[0,n]}(x)$.
(c) $f_n(x)=e^{-x/n}\,e^{-x}$.
Justify each answer by naming a dominating function or an increasing structure
where one exists.
\end{exercise}

\begin{exercise}[Mixed random variable]
\label{ch37:ex-mixed}
Let $X=0$ with probability $p$ and $X\sim\Unif(0,1)$ with probability $1-p$.
(a) Write $\E[X]$ as a sum of two Lebesgue integrals over complementary events.
(b) Compute $\E[X]$ and $\E[X^2]$, and hence $\Var(X)$, in terms of $p$.
(c) Explain why neither the pure discrete nor the pure continuous expectation
formula computes this directly, but the Lebesgue definition does.
\end{exercise}

\begin{exercise}[Term-by-term integration]
\label{ch37:ex-termbyterm}
Let $g_k(x)=x^{k}(1-x)$ on $[0,1]$ for $k=0,1,2,\dots$ (all nonnegative).
(a) Use MCT (via partial sums) to justify $\int_0^1 \sum_{k\ge0} g_k\,\dd x
= \sum_{k\ge0}\int_0^1 g_k\,\dd x$. (b) Evaluate both sides
($\sum_{k\ge0} g_k(x)=1$ for $x\in[0,1)$) and confirm they agree. (c) Where,
exactly, did nonnegativity of the $g_k$ matter?
\end{exercise}

\begin{exercise}[Riemann inside Lebesgue]
\label{ch37:ex-riemann}
(a) State the criterion from Theorem~\ref{ch37:thm-riemann} for a bounded
function on $[a,b]$ to be Riemann integrable. (b) Use it to decide whether
$f(x)=\ind_{\{x \text{ is a dyadic rational}\}}$ on $[0,1]$ is Riemann
integrable, and compute its Lebesgue integral. (c) Give a bounded function on
$[0,1]$ that is Lebesgue integrable but not Riemann integrable, and one
(improper) example that is Riemann integrable but not Lebesgue integrable,
explaining the difference in one sentence.
\end{exercise}

\begin{exercise}[Differentiating under the integral, conceptual]
\label{ch37:ex-diffunder}
Suppose $X$ has density $f$ and $g(x,\theta)$ is differentiable in $\theta$ with
$\abs{\partial_\theta g(x,\theta)}\le h(x)$ where $\E[h(X)]<\infty$. Explain,
citing the relevant convergence theorem, why
$\frac{\dd}{\dd\theta}\E[g(X,\theta)]=\E[\partial_\theta g(X,\theta)]$. (Hint:
write the derivative as a limit of difference quotients and dominate them with
$h$ via the mean value theorem.) Why is this the workhorse behind score
functions and Fisher information in maximum-likelihood estimation?
\end{exercise}
